\label{sect:environments}
\label{sect:examples-venv}

The following documentation is specific to the \textbf{Speed} HPC Facility at the
Gina Cody School of Engineering and Computer Science.
%
Virtual environments typically instantiated via Conda or Python.
Another option is Singularity detailed in \xs{sect:singularity-containers}.
%
Usually, virtual environments are created once during an interactive session before submitting
a batch job to the scheduler. The job script submitted to the scheduler is then written to 
(1) activate the virtual environment, (2) use it, and (3) close it at the end of the job.

% ------------------------------------------------------------------------------
\subsubsection{Anaconda}
\label{sect:conda-venv}

Request an interactive session in the queue you wish to submit your jobs to (e.g., salloc -p pg --gpus=1 for GPU jobs).
Once your interactive has started, create an anaconda environment in your speed-scratch directory by using 
the \texttt{\-\-prefix} option when executing \texttt{conda create}. For example, to create an anaconda environment for 
\texttt{a\_user}, execute the following at the command line:

\begin{verbatim}
module load anaconda3/2023.03/default
conda create --prefix /speed-scratch/a_user/myconda
\end{verbatim}

\vspace{10pt}
\noindent
\textbf{Note:} Without the \texttt{\-\-prefix} option, the \texttt{conda create} command creates the 
environment in \texttt{a\_user}'s home directory by default.
\vspace{10pt}

% ------------------------------------------------------------------------------
\paragraph{List Environments.}

To view your conda environments, type: \texttt{conda info --envs}

\begin{verbatim}
# conda environments:
#
base                  *  /encs/pkg/anaconda3-2023.03/root
                         /speed-scratch/a_user/myconda
\end{verbatim}      

% ------------------------------------------------------------------------------
\paragraph{Activate an Environment.}

Activate the environment \texttt{\/speed\-scratch\/a\_user\/myconda} as follows
\begin{verbatim}
conda activate /speed-scratch/a_user/myconda
\end{verbatim}
After activating your environment, add \tool{pip} to your environment by using 
\begin{verbatim}
conda install pip
\end{verbatim}
This will install \tool{pip} and \tool{pip}'s dependencies, including python, 
into the environment.

\begin{itemize}
	\item
		A consolidated example using Conda:
		\small
		\begin{verbatim}
			salloc -p pg --gpus=1 --mem=10GB -A <slurm account name>
			cd /speed-scratch/$USER
			module load python/3.11.0/default
			conda create -p /speed-scratch/$USER/pytorch-env
			conda activate /speed-scratch/$USER/pytorch-env
			conda install python=3.11.0
			pip3 install torch torchvision torchaudio --index-url \ 
			  https://download.pytorch.org/whl/cu117
			....
			conda deactivate
			exit # end the salloc session
		\end{verbatim}
		\normalsize
	\item
	    No Space left error: 
			\texttt Read our Github
		\href
		{https://github.com/NAG-DevOps/speed-hpc/tree/master/src#no-space-left-error-when-creating-conda-environment}
		{\texttt{HERE}}
\end{itemize}

\noindent
\textbf{Important Note:} \tool{pip} (and \tool{pip3}) are used to install modules
from the python distribution while \texttt{conda install} installs modules from 
anaconda's repository.

% ------------------------------------------------------------------------------
\paragraph{Conda Env without --prefix: }

If you don't want to use the \texttt{\-\-prefix} option every time you create a new environment and you don't want to use the default \texttt{\-\$HOME}.
Create a new directory an set the following variables to point to the new created directory, e.g:
\begin{verbatim}
	setenv CONDA_ENVS_PATH /speed-scratch/$USER/condas
	setenv CONDA_PKGS_DIRS /speed-scratch/$USER/condas/pkg
\end{verbatim}
If you want to make these changes permanent, add the variables to your \texttt{.tcshrc} or \texttt{.bashrc} (depending on the default shell you are using)

% ------------------------------------------------------------------------------
\subsubsection{Python}
\label{sect:python-venv}

Setting up a Python virtual environment is fairly straightforward.
The first step is to request an interactive session in the queue you wish to submit your jobs to.

We have a simple example that use a Python virtual environment:

\begin{itemize}
	\item
		Using Python Venv
		\small
		\begin{verbatim}
			salloc -p pg --gpus=1 --mem=10GB -A <slurm account name>
			cd /speed-scratch/$USER
			module load python/3.9.1/default
			mkdir -p /speed-scratch/$USER/tmp 
			setenv TMPDIR /speed-scratch/$USER/tmp
			setenv TMP /speed-scratch/$USER/tmp
			python -m venv $TMPDIR/testenv (testenv=name of the virtualEnv)
			source /speed-scratch/$USER/tmp/testenv/bin/activate.csh
			pip install modules…
			deactivate
			exit
		\end{verbatim}
		\normalsize
	\item
	    See, e.g., 
		\href
		{https://github.com/NAG-DevOps/speed-hpc/blob/master/src/gurobi-with-python.sh}
		{\texttt{gurobi-with-python.sh}}
\end{itemize}

\noindent
\textbf{Important Note:} partition \texttt{ps} is used for CPU jobs, partitions \texttt{pg}, \texttt{pt} are used
for GPU jobs, no need to use \texttt{--gpus=} when preparing environments for CPU jobs.